% sec-01: why the size is exactly $2,104,000, and the seven proceeds lines.
% Table 6.  No figure: the arithmetic is four numbers and the allocation is a
% table, and a figure of either would restate rather than add.
\section{Why the Size Is \$2,104,000}

A private raise sized at anything other than the gap is either asking for money
the federal route already supplies or leaving a hole that a later round has to
fill at a worse price. The gap is four numbers, none of them derived here.

The five-year program is \$3,500,000 in direct cost, at \$700,000 per year, and
that frame is reused from the ten-application work rather than recalculated
\cite{tenapps2026}. The Small Business Innovation Research route is \$306,000
over nine months in Phase I and \$1,300,000 over twenty-four months in Phase II,
which is \$1,606,000 over thirty-three months \cite{nihsbirseed,sbapolicydirective}.
Of that total, \$1,396,000 is direct work inside the award once the fee and the
indirect allowance are set aside \cite{ecfr2026uniformguidance,nihgpsfee}. The
program's direct cost less that figure is \$2,104,000 \cite{capitalplan2026}.

\subsection{The seven lines the proceeds would fund}

\begin{apptable}
\begin{tabularx}{\textwidth}{>{\raggedright\arraybackslash}p{4.4cm} >{\raggedright\arraybackslash}p{2.0cm} >{\raggedright\arraybackslash}X}
\headrow \thead{Line} & \thead{Amount} & \thead{What it buys that the federal route does not}\\
\midrule
Site start-up and feasibility & \$420,000 & Site agreement, coverage analysis, and the institutional review the award funds only after it is made \\
Clinical operations beyond the award & \$560,000 & Coordinator, regulatory and quality time in the months the award does not cover \\
Robotic platform access and integration & \$380,000 & Platform time, integration testing, and the bench evidence a site requires before it agrees \\
Regulatory submissions and meetings & \$215,000 & A Pre-Request for Designation, a pre-IND Type B meeting, and a device Q-Submission \\
Insurance, indemnification, subject injury reserve & \$260,000 & The cover no site proceeds without and no federal award pays for directly \\
Verification and model governance beyond Phase I & \$180,000 & The verification harness maintained past the nine-month Phase I window \\
Working capital and contingency & \$89,000 & The residual, held rather than allocated \\
\midrule
\textbf{Total} & \textbf{\$2,104,000} & Equal to the gap above, by construction \\
\end{tabularx}
\end{apptable}
\vspace{-0.60cm}
\tabcap{Table 6. The seven proceeds lines, each with what it buys that a federal\\
award does not, summing to the gap the federal route leaves open.}

\subsection{What is never paid from federal award funds}

A private holder is entitled to know that their money and federal money are not
fungible in this company. The separation below is the one the site documentation
package sets out, and the rule is not a preference: the same activity can be
lawful with one source of funds and unlawful with another \cite{movein2026}.

Lobbying of any kind, the cost of raising this financing, entertainment, any cost
outside the approved award scope, and any cost incurred before an award start
date are paid from company or private funds and never from a federal award. The
reverse also holds: work inside an approved award scope is charged to the award
and not to private capital, so that a holder is not quietly subsidizing something
the government has already paid for.

\subsection{The cost of publishing this}

A company that publishes its capitalization detail gives up three things, and the
capitalization plan says so in its own risk section rather than hiding it
\cite{capitalplan2026}. The individual-scientist framing that made the original
applications distinctive is set aside. Detail an owner might prefer private
becomes public. And some prior recipients will read a company framing as
off-mechanism for the person-based awards they administer.

That trade was made deliberately. It is restated here so that a reader who finds
the capitalization detail in a public repository knows the disclosure was a
choice rather than an accident.

\subsection{What a reader should check before believing any of this}

Every figure above is checkable against the public repository without contacting
this company, which is the point of publishing it. Four checks in particular are
worth five minutes each.

That the \$3,500,000 five-year and \$700,000 annual figures appear in the
ten-application work at the same values, unamended \cite{tenapps2026}. That the
\$306,000 and \$1,300,000 Small Business Innovation Research split appears in the
application addressed to that mechanism, and that the phase durations match
\cite{nihsbirseed}. That every quantitative result in the evidence pack carries
the limitation its own authors stated, in the same row as the result
\cite{capitalplan2026}. And that no document anywhere in the repository describes
any institution as a partner, sponsor, site, or endorser \cite{ucsdmoores2026}.

A reader who runs those four checks and finds a discrepancy has found something
this company would want to know about, and the repository's issue tracker is the
fastest route to telling it.

There is a fifth check that cannot be run against the repository at all, and it
is named here because its absence is the honest limit of everything above. No
document in this repository can establish that a clinical institution will take
the study, because no institution has been asked in a form that would produce a
binding answer, and the letters this day sends are first approaches rather than
proposals. A reader who wants that assurance has to wait for the institutions
themselves to answer, and this company will publish those answers whichever way
they fall.

